%! Describing the Distribution of a Quantitative Variable — Unit 1, Lesson 1.4 — Warm-Up
% ONE PAGE, blank and key. Three items, set at 12pt like every other student component.
% Do not add a fourth item — the page is full.
% GRADUAL RELEASE: the warm-up activates prior knowledge before the guided notes. Item 1 is
% the spiral back to Lesson 1.3 (two interval widths, one data set); item 2 is the
% count-to-the-middle move notes section 1 needs, so finding a center is not the obstacle
% there; item 3 puts a lone far-off value in front of them in plain language — notes section 1
% names it an OUTLIER ten minutes later, and section 2 names the lean.
% NO SKETCH QUESTIONS: the item 3 display is pre-drawn; students only read it.
\documentclass[12pt]{article}
\usepackage{apstats-article}
\usepackage{apstats-boxes}

\begin{document}

\pageheader{Unit 1, Lesson 1.4}{Warm-Up}

\vspace{0.20in}

\begin{enumerate}[leftmargin=1.9em, itemsep=0.42in]

  \item \textbf{From last lesson.} One class's 40 quiz scores were drawn twice. With blocks of
        5 points the picture showed two separate peaks; with blocks of 20 points the very same
        40 scores showed one single hump.

        \vspace{10pt}
        Which of the two pictures is \emph{wrong}? \blank{3.4cm}

        \vspace{14pt}
        What should you check before you believe what any histogram looks like?
        \blank{4.6cm}

  \item Nine customers were timed at a coffee counter. Their waiting times, in minutes, are
        already written in order from shortest to longest.

        \vspace{10pt}
        \begin{center}
        1 \quad 2 \quad 2 \quad 3 \quad 4 \quad 4 \quad 5 \quad 7 \quad 12
        \end{center}
        \vspace{10pt}

        How many waits are listed? \blank{1.6cm} \qquad
        Counting in from both ends, which value lands exactly in the middle? \blank{1.6cm}

  \item Fifteen students were timed waiting for a ride home after school.

        \vspace{4pt}
        \begin{center}
        \begin{tikzpicture}[scale=0.9]
          \draw[->, thick] (-0.2,0) -- (8.6,0);
          \foreach \x/\lab in {0/0, 1/2, 2/4, 3/6, 4/8, 5/10, 6/12, 7/14, 8/16} {
            \draw (\x,-0.07) -- (\x,0.07); \node[below, font=\tiny] at (\x,-0.1) {\lab};}
          \foreach \x/\n in {1/2, 1.5/3, 2/4, 2.5/2, 3/2, 3.5/1, 7.5/1} {
            \foreach \k in {1,...,\n} { \fill[navy] (\x,{0.22*\k}) circle (0.075); }}
          \node[font=\tiny] at (4.2,-0.85) {Minutes waited for a ride};
        \end{tikzpicture}
        \end{center}
        \vspace{4pt}

        Most of the waits sit between \blank{1.3cm} and \blank{1.3cm} minutes. \quad
        One wait sits far away from all the others: \blank{1.6cm} minutes.

        \vspace{10pt}
        In one sentence, describe where the other fourteen waits are compared with that one.\\[10pt]
        \writeline\\[14pt]\writeline

\end{enumerate}

\end{document}
